% chapters/06/section-05.tex
\section{Program Fisika Lengkap}
\label{sec:program-fisika}

Bagian ini menyatukan ketiga unsur fisika menjadi satu program yang memuat dua
benda dengan perilaku berbeda, sekaligus mengukur hasilnya untuk dibandingkan
dengan perhitungan.

\subsection{Kode Lengkap}

\kodepython{python/ch06_fisika.py}{Gravitasi, lompatan, dan pantulan pada dua benda}{lst:ch06-fisika}

\subsection{Dua Benda, Satu Hukum}

Program tersebut memuat dua benda yang perilakunya berbeda, tetapi dihitung
oleh fungsi yang sama.

\begin{description}
  \item[Pemain] digerakkan tombol mendatar, jatuh karena gravitasi, berhenti di
        lantai karena koefisien restitusinya nol, dan dapat melompat ketika
        menyentuh lantai. Benda ini memperlihatkan indikator \ind{a} dan
        \ind{b}.
  \item[Bola] bergerak sendiri, memantul di lantai dan di kedua dinding dengan
        koefisien restitusi 0,8 sehingga pantulannya makin lama makin pendek.
        Benda ini memperlihatkan indikator \ind{c}.
\end{description}

Perbedaan perilaku keduanya hanya berasal dari dua hal: siapa
yang menentukan kecepatan mendatarnya, dan berapa nilai koefisien restitusinya.
Seluruh perhitungan lainnya, yaitu gravitasi, lantai, dan dinding, sama persis.

\subsection{Mengukur Tinggi Lompatan dan Membandingkannya dengan Hitungan}

Program bab ini mencatat titik tertinggi yang dicapai pemain, lalu
menampilkannya bersama tinggi yang seharusnya menurut Persamaan~\eqref{eq:tinggi-lompat}.
Perbandingan keduanya pada beberapa laju bingkai diperlihatkan pada
Tabel~\ref{tab:galat-euler}.

\begin{table}[htbp]
\centering
\caption{Tinggi lompatan terukur dibandingkan hitungan, pada beberapa laju bingkai}
\label{tab:galat-euler}
{\setlength{\tabcolsep}{4pt}\footnotesize%
\begin{tabularx}{\textwidth}{@{}C{1.8cm} C{2.6cm} C{2.9cm} Y@{}}
\toprule
\headrow \textbf{Batas laju} & \textbf{Tinggi terukur} & \textbf{Menurut perhitungan} & \textbf{Galat} \\
\midrule
30 FPS & 124,6 piksel & 124,7 piksel & 11,4 piksel (8,4\%) \\
\addlinespace
60 FPS & 130,3 piksel & 130,3 piksel & 5,8 piksel (4,3\%) \\
\addlinespace
120 FPS & 133,1 piksel & 133,2 piksel & 2,9 piksel (2,1\%) \\
\addlinespace
240 FPS & 134,5 piksel & 134,7 piksel & 1,5 piksel (1,1\%) \\
\bottomrule
\end{tabularx}}
\tabsource{Gravitasi 1800 piksel per detik kuadrat dan kecepatan lompatan 700
piksel per detik, sehingga tinggi yang seharusnya adalah 136,1 piksel. Kolom
ketiga dihitung dengan menjumlahkan perpindahan setiap bingkai.}
\end{table}

Tiga hal dapat dibaca dari tabel tersebut.

\begin{enumerate}
  \item \textbf{Hasil pengukuran sesuai dengan perhitungan.} Tinggi terukur
        dan hasil perhitungan berbeda kurang dari 0,2 piksel pada setiap laju.
        Kesesuaian itu menunjukkan bahwa simulasi berjalan sebagaimana
        dimaksudkan, dan bahwa cara pengukurannya benar.
  \item \textbf{Simulasi selalu lebih rendah daripada lintasan sebenarnya.}
        Tinggi yang seharusnya adalah 136,1 piksel, sedangkan simulasi hanya
        mencapai 124 sampai 135 piksel. Selisih itu adalah galat integrasi
        Euler yang telah dijelaskan pada Bagian~2.4, bukan kesalahan penulisan
        program.
  \item \textbf{Galat menyusut separuh setiap panjang langkah diperkecil
        separuh.} Perhatikan kolom terakhir: 11,4 lalu 5,8 lalu 2,9 lalu 1,5.
        Pola itu sesuai dengan sifat ketelitian tingkat pertama yang disebutkan
        pada Bagian~2.4.
\end{enumerate}

Pelajaran praktisnya penting bagi perancangan permainan. Pada 60 bingkai per
detik, galat sebesar 4,3 persen tidak dapat dirasakan pemain, karena tingginya
hanya berbeda enam piksel pada lompatan setinggi 136 piksel. Karena itu, kaidah
Euler yang sederhana sudah memadai untuk sebagian besar permainan, dan
perhitungan yang lebih rumit hanya diperlukan apabila hasil simulasi harus
dapat diulang secara persis.

\begin{catatabox}
Lompatan pemain sebenarnya \textbf{tidak persis} setinggi yang
dihitung dengan Persamaan~\eqref{eq:tinggi-lompat}. Apabila permainan Anda
menuntut pemain mencapai bongkahan yang berada tepat di batas jangkauan
lompatan, sediakan sedikit kelonggaran. Kebiasaan memberikan kelonggaran
sepuluh sampai dua puluh persen merupakan kebiasaan yang lazim pada
perancangan level.
\end{catatabox}

\subsection{Percobaan yang Disarankan}

\begin{table}[htbp]
\centering
\caption{Percobaan pada program fisika}
\label{tab:percobaan-ch06}
{\setlength{\tabcolsep}{4pt}\footnotesize%
\begin{tabularx}{\textwidth}{@{}L{5.6cm} Y@{}}
\toprule
\headrow \textbf{Yang diubah dan caranya} & \textbf{Yang sebaiknya Anda amati} \\
\midrule
Gravitasi, melalui pilihan \code{--gravity}. &
Lompatan menjadi lebih tinggi dan lebih melayang ketika gravitasi diperkecil,
walaupun kecepatan lompatannya tidak berubah. \\
\addlinespace
Kecepatan lompatan, melalui pilihan \code{--jump-speed}. &
Tinggi lompatan berubah sebanding dengan kuadrat kecepatan, bukan sebanding
langsung. Menaikkan kecepatan dua kali menaikkan tinggi empat kali. \\
\addlinespace
Koefisien restitusi bola, melalui pilihan \code{--restitution}. &
Pada nilai 1, bola memantul selamanya dengan tinggi yang tetap. Pada nilai 0,
bola berhenti seketika seperti pemain. \\
\addlinespace
Koefisien restitusi pemain, dengan mengubah \code{PLAYER_RESTITUTION}. &
Pemain memantul setiap kali mendarat dan kendalinya terasa lepas. Kembalikan
nilainya menjadi nol sesudah percobaan. \\
\addlinespace
Cara membaca lompatan, dengan mengganti \code{KEYDOWN} menjadi
\code{key.get_pressed()}. &
Menahan tombol lompatan membuat pemain melompat terus-menerus. Kembalikan
sesudah percobaan. \\
\bottomrule
\end{tabularx}}
\end{table}
